% Options for packages loaded elsewhere
\PassOptionsToPackage{unicode}{hyperref}
\PassOptionsToPackage{hyphens}{url}
\PassOptionsToPackage{dvipsnames,svgnames,x11names}{xcolor}
%
\documentclass[
  11pt,
]{article}
\usepackage{amsmath,amssymb}
\usepackage{iftex}
\ifPDFTeX
  \usepackage[T1]{fontenc}
  \usepackage[utf8]{inputenc}
  \usepackage{textcomp} % provide euro and other symbols
\else % if luatex or xetex
  \usepackage{unicode-math} % this also loads fontspec
  \defaultfontfeatures{Scale=MatchLowercase}
  \defaultfontfeatures[\rmfamily]{Ligatures=TeX,Scale=1}
\fi
\usepackage{lmodern}
\ifPDFTeX\else
  % xetex/luatex font selection
\fi
% Use upquote if available, for straight quotes in verbatim environments
\IfFileExists{upquote.sty}{\usepackage{upquote}}{}
\IfFileExists{microtype.sty}{% use microtype if available
  \usepackage[]{microtype}
  \UseMicrotypeSet[protrusion]{basicmath} % disable protrusion for tt fonts
}{}
\makeatletter
\@ifundefined{KOMAClassName}{% if non-KOMA class
  \IfFileExists{parskip.sty}{%
    \usepackage{parskip}
  }{% else
    \setlength{\parindent}{0pt}
    \setlength{\parskip}{6pt plus 2pt minus 1pt}}
}{% if KOMA class
  \KOMAoptions{parskip=half}}
\makeatother
\usepackage{xcolor}
\usepackage[margin=1in]{geometry}
\usepackage{longtable,booktabs,array}
\usepackage{calc} % for calculating minipage widths
% Correct order of tables after \paragraph or \subparagraph
\usepackage{etoolbox}
\makeatletter
\patchcmd\longtable{\par}{\if@noskipsec\mbox{}\fi\par}{}{}
\makeatother
% Allow footnotes in longtable head/foot
\IfFileExists{footnotehyper.sty}{\usepackage{footnotehyper}}{\usepackage{footnote}}
\makesavenoteenv{longtable}
\setlength{\emergencystretch}{3em} % prevent overfull lines
\providecommand{\tightlist}{%
  \setlength{\itemsep}{0pt}\setlength{\parskip}{0pt}}
\setcounter{secnumdepth}{-\maxdimen} % remove section numbering
\ifLuaTeX
  \usepackage{selnolig}  % disable illegal ligatures
\fi
\IfFileExists{bookmark.sty}{\usepackage{bookmark}}{\usepackage{hyperref}}
\IfFileExists{xurl.sty}{\usepackage{xurl}}{} % add URL line breaks if available
\urlstyle{same}
\hypersetup{
  pdfauthor={Stefano Alimonti --- stefalimonti@gmail.com},
  colorlinks=true,
  linkcolor={blue},
  filecolor={Maroon},
  citecolor={Blue},
  urlcolor={blue},
  pdfcreator={LaTeX via pandoc}}

\author{Stefano Alimonti --- stefalimonti@gmail.com}
\date{March 2026}

\begin{document}

{
\hypersetup{linkcolor=black}
\setcounter{tocdepth}{2}
\tableofcontents
}
\section{Exact Covariance Structure of Binary Carry
Chains}\label{exact-covariance-structure-of-binary-carry-chains}

\textbf{Author:} Stefano Alimonti \textbf{Affiliation:} Independent Researcher
\textbf{Date:} March 2026

\begin{center}\rule{0.5\linewidth}{0.5pt}\end{center}

\textbf{Abstract.} We establish the exact second-order structure of the carry
chain arising from binary multiplication with fixed most-significant bits. Six
results are proved: (i) the column expectation E{[}conv\_j{]} = (j+3)/4 (Lemma
A); (ii) the Parity Lemma, that total\_j mod 2 is uniform for all j ≥ 1 (Lemma
B); (iii) the carry expectation E{[}carry\_j{]} = (j−1)/4 (Theorem 3); (iv) the
universal off-diagonal covariance Cov(carry\_j, g\_i h\_\{j−i\}) = 1/8 for all
off-diagonal pairs (Lemma D); (v) Cov(carry\_j, conv\_j) = (j−1)/8 for all odd j
≥ 3 (Main Theorem); and (vi) the exact inductive step ΔCov = 1/4 for odd-to-odd
transitions (Corollary). For even j, a diagonal correction ε\_j \textgreater{} 0
arises from the failure of the Parity Lemma under conditioning; we tabulate
these corrections through j = 12.

\textbf{Overview.} Section 1 introduces the binary multiplication model. Section
2 establishes three fundamental identities: the column expectation (Lemma A),
the Parity Lemma (Lemma B), and the carry expectation (Theorem 3). Section 3
proves the universal off-diagonal covariance (Lemma D). Section 4 derives the
main covariance formula and its corollary. Section 5 treats even-j corrections.
Sections 6--7 present computational verification and discussion.

\begin{center}\rule{0.5\linewidth}{0.5pt}\end{center}

\newpage

\subsection{1. Setup and Notation}\label{setup-and-notation}

We study the carry chain arising from binary multiplication of two n-bit
integers G = g₀ g₁ ⋯ g\_\{n−1\} and H = h₀ h₁ ⋯ h\_\{n−1\} in base 2.

\textbf{Bit model.} Fix the most-significant bits g₀ = h₀ = 1. For i ≥ 1, each
g\_i and h\_i is drawn independently from Bernoulli(1/2).

\textbf{Convolution sums.} For each column j ≥ 0:

\[\text{conv}_j = \sum_{i=0}^{j} g_i \, h_{j-i}.\]

\textbf{Carry recurrence.} Set carry₀ = 0, and define

\[\text{carry}_{j+1} = \left\lfloor \frac{\text{conv}_j + \text{carry}_j}{2} \right\rfloor.\]

Write total\_j = conv\_j + carry\_j for the full column sum. The output bit at
position j is total\_j mod 2, and the carry propagated forward is ⌊total\_j /
2⌋.

\textbf{Terminology.} A bit-product term g\_i h\_\{j−i\} appearing in conv\_j is
\emph{off-diagonal} if i ≠ j − i (equivalently, j is odd, or j is even and i ≠
j/2). It is \emph{diagonal} if i = j − i = j/2.

\begin{center}\rule{0.5\linewidth}{0.5pt}\end{center}

\newpage

\subsection{2. Fundamental Identities}\label{fundamental-identities}

\subsubsection{2.1 Lemma A (Column
Expectation)}\label{lemma-a-column-expectation}

\textbf{Lemma A.} \emph{For j ≥ 1,}

\[E[\text{conv}_j] = \frac{j+3}{4}.\]

\emph{Proof.} The convolution sum conv\_j contains j + 1 terms. Two are border
terms involving the fixed MSBs: g₀ h\_j = h\_j with E{[}h\_j{]} = 1/2, and g\_j
h₀ = g\_j with E{[}g\_j{]} = 1/2. The remaining j − 1 terms are interior
products g\_i h\_\{j−i\} with both factors random, so E{[}g\_i h\_\{j−i\}{]} =
1/4. Therefore:

\[E[\text{conv}_j] = 2 \cdot \frac{1}{2} + (j-1) \cdot \frac{1}{4} = 1 + \frac{j-1}{4} = \frac{j+3}{4}. \qquad \square\]

\subsubsection{2.2 Lemma B (Parity Lemma)}\label{lemma-b-parity-lemma}

\textbf{Lemma B} (Parity Lemma). \emph{For all j ≥ 1,}

\[P(\text{total}_j \text{ is odd}) = \frac{1}{2}.\]

\emph{Proof.} Write total\_j = g\_j · h₀ + W\_j, where

\[W_j = g_0 \cdot h_j + \sum_{i=1}^{j-1} g_i h_{j-i} + \text{carry}_j.\]

Since h₀ = 1, we have total\_j = g\_j + W\_j. The digit g\_j first enters the
computation at column j: it does not appear in conv₀, \ldots, conv\_\{j−1\} and
hence is independent of carry\_j. It is also independent of every term in W\_j.
Since g\_j ∼ Ber(1/2):

\[P(\text{total}_j \text{ odd}) = \tfrac{1}{2}\,P(W_j \text{ even}) + \tfrac{1}{2}\,P(W_j \text{ odd}) = \frac{1}{2}. \qquad \square\]

\textbf{Remark} (General base). For base b ≥ 2, the analogous statement is that
total\_j mod b is uniform on \{0, \ldots, b−1\}. The proof uses the same
decomposition total\_j = g\_j · h₀ + W\_j: as g\_j ranges over \{0, \ldots,
b−1\}, the map a ↦ a · h₀ mod b is a bijection if and only if gcd(h₀, b) = 1.
For base 2 this is automatic since h₀ = 1. For general bases, the coprimality
condition gcd(g₀ h₀, b) = 1 suffices, since either g\_j · h₀ or g₀ · h\_j
provides the required uniform residue. When both gcd(g₀, b) \textgreater{} 1 and
gcd(h₀, b) \textgreater{} 1, a character-sum argument using the pair (g\_j,
h\_j) simultaneously shows that uniform distribution still holds provided
gcd(g₀, b) · gcd(h₀, b) \textless{} b, which is satisfied for all valid MSB
pairs in bases b ≤ 12 and for all prime-power bases.

\subsubsection{2.3 Theorem 3 (Carry
Expectation)}\label{theorem-3-carry-expectation}

\textbf{Theorem 3.} \emph{For all j ≥ 1,}

\[E[\text{carry}_j] = \frac{j-1}{4}.\]

\emph{Proof.} By induction. \textbf{Base:} carry₁ = ⌊g₀ h₀ / 2⌋ = ⌊1/2⌋ = 0 = (1
− 1)/4. ✓

\textbf{Inductive step.} Suppose E{[}carry\_j{]} = (j − 1)/4. Since
carry\_\{j+1\} = (total\_j − (total\_j mod 2)) / 2:

\[E[\text{carry}_{j+1}] = \frac{E[\text{total}_j] - E[\text{total}_j \bmod 2]}{2}.\]

By Lemma A and the inductive hypothesis, E{[}total\_j{]} = E{[}conv\_j{]} +
E{[}carry\_j{]} = (j + 3)/4 + (j − 1)/4 = (j + 1)/2. By the Parity Lemma,
E{[}total\_j mod 2{]} = P(total\_j odd) = 1/2. Hence:

\[E[\text{carry}_{j+1}] = \frac{(j+1)/2 - 1/2}{2} = \frac{j}{4}. \qquad \square\]

\begin{center}\rule{0.5\linewidth}{0.5pt}\end{center}

\newpage

\subsection{3. Universal Off-Diagonal Covariance (Lemma
D)}\label{universal-off-diagonal-covariance-lemma-d}

\textbf{Lemma D.} \emph{For all 1 ≤ i ≤ j − 1 with i ≠ j − i:}

\[\text{Cov}(\text{carry}_j,\; g_i h_{j-i}) = \frac{1}{8}.\]

\emph{Proof.} Since g\_i h\_\{j−i\} ∈ \{0, 1\} with P(g\_i h\_\{j−i\} = 1) =
1/4:

\[\text{Cov}(\text{carry}_j,\; g_i h_{j-i}) = \frac{1}{4}\,\delta E[c_j],\]

where δE{[}c\_j{]} = E{[}carry\_j \textbar{} g\_i = 1, h\_\{j−i\} = 1{]} −
E{[}carry\_j{]} is the perturbation in the carry expectation caused by
conditioning. We establish δE{[}c\_j{]} = 1/2 in four steps.

\textbf{Step 1: The Parity Lemma survives conditioning.} Let a = min(i, j − i)
and b = max(i, j − i). Since i ≠ j − i, we have a \textless{} b, and the
conditioning fixes at most one bit at any given index. Specifically: at column
k, the fresh bit g\_k is conditioned only if k = i, but then h\_k is free (since
k = i ≠ j − i means h\_k is unconditioned), and h\_k can serve as the uniform
parity-flipping variable. Similarly, if k = j − i, then g\_k is free. At all
other columns, g\_k is free. Therefore the Parity Lemma holds at every column k
≥ 1 under the conditioning.

\textbf{Step 2: Perturbation recursion.} Let δE{[}c\_k{]} = E{[}carry\_k
\textbar{} g\_i = 1, h\_\{j−i\} = 1{]} − E{[}carry\_k{]}. Because the Parity
Lemma holds under conditioning, the mod-2 correction is unchanged, giving the
exact recursion:

\[\delta E[c_{k+1}] = \frac{\delta E[\text{conv}_k] + \delta E[c_k]}{2},\]

with δE{[}c₀{]} = 0.

\textbf{Step 3: Perturbation profile δE{[}conv\_k{]}.} Conditioning g\_i = 1
shifts E{[}g\_i{]} from 1/2 to 1; conditioning h\_\{j−i\} = 1 shifts
E{[}h\_\{j−i\}{]} from 1/2 to 1. The term g\_i h\_\{k−i\} in conv\_k gains +1/2
(border, at k = i) or +1/4 (interior, at k \textgreater{} i); similarly for
h\_\{j−i\}. The full profile is:

\begin{longtable}[]{@{}ll@{}}
\toprule\noalign{}
Range & δE{[}conv\_k{]} \\
\midrule\noalign{}
\endhead
\bottomrule\noalign{}
\endlastfoot
k \textless{} a & 0 \\
k = a & 1/2 \\
a \textless{} k \textless{} b & 1/4 \\
k = b & 3/4 \\
k \textgreater{} b & 1/2 \\
\end{longtable}

At k = a, the first conditioned bit enters as a border term (+1/2). For a
\textless{} k \textless{} b, it contributes as an interior term (+1/4). At k =
b, the second conditioned bit enters as a border term (+1/2) while the first
continues as interior (+1/4), giving 3/4. For k \textgreater{} b, both
contribute interior terms (+1/4 each).

\textbf{Step 4: Exact evaluation via generating function.} Define T\_k = 2\^{}k
· δE{[}c\_k{]}. The recursion yields T\_\{k+1\} = T\_k + 2\^{}k ·
δE{[}conv\_k{]}, so:

\[T_j = \sum_{k=0}^{j-1} 2^k \,\delta E[\text{conv}_k].\]

Substituting the profile from Step 3:

\[T_j = 2^a \cdot \frac{1}{2} + \frac{1}{4}\sum_{k=a+1}^{b-1} 2^k + \frac{3}{4} \cdot 2^b + \frac{1}{2}\sum_{k=b+1}^{j-1} 2^k.\]

Evaluating the geometric sums:

\[T_j = 2^{a-1} + \frac{2^b - 2^{a+1}}{4} + \frac{3 \cdot 2^b}{4} + \frac{2^j - 2^{b+1}}{2}.\]

Collecting terms:

\[T_j = 2^{a-1} + \frac{4 \cdot 2^b - 2^{a+1}}{4} + \frac{2^j - 2^{b+1}}{2} = 2^{a-1} + 2^b - 2^{a-1} + 2^{j-1} - 2^b = 2^{j-1}.\]

Therefore δE{[}c\_j{]} = T\_j / 2\^{}j = 1/2, and:

\[\text{Cov}(\text{carry}_j,\; g_i h_{j-i}) = \frac{1}{4} \cdot \frac{1}{2} = \frac{1}{8}. \qquad \square\]

\begin{center}\rule{0.5\linewidth}{0.5pt}\end{center}

\newpage

\subsection{4. The Main Covariance Theorem}\label{the-main-covariance-theorem}

\textbf{Main Theorem.} \emph{For all odd j ≥ 3:}

\[\text{Cov}(\text{carry}_j, \text{conv}_j) = \frac{j-1}{8}.\]

\emph{Proof.} Expand by linearity of covariance:

\[\text{Cov}(\text{carry}_j, \text{conv}_j) = \sum_{i=0}^{j} \text{Cov}(\text{carry}_j,\; g_i h_{j-i}).\]

We classify each term by type:

\begin{itemize}
\item
  \textbf{Border terms (i = 0 and i = j):} Cov(carry\_j, g₀ h\_j) =
  Cov(carry\_j, h\_j) = 0, because h\_j first enters the computation at column j
  and is independent of carry\_j (which depends only on columns 0, \ldots, j −
  1). Similarly, Cov(carry\_j, g\_j h₀) = Cov(carry\_j, g\_j) = 0.
\item
  \textbf{Off-diagonal terms (1 ≤ i ≤ j − 1):} Since j is odd, i ≠ j − i for
  every such i (there is no diagonal term). By Lemma D, each covariance equals
  1/8.
\end{itemize}

There are exactly j − 1 off-diagonal terms, giving:

\[\text{Cov}(\text{carry}_j, \text{conv}_j) = (j-1) \cdot \frac{1}{8} = \frac{j-1}{8}. \qquad \square\]

\textbf{Corollary} (Exact inductive step). \emph{For all odd j ≥ 3:}

\[\text{Cov}(\text{carry}_{j+2}, \text{conv}_{j+2}) - \text{Cov}(\text{carry}_j, \text{conv}_j) = \frac{1}{4}.\]

\emph{Proof.} Direct computation:

\[\frac{(j+2) - 1}{8} - \frac{j - 1}{8} = \frac{j+1 - j + 1}{8} = \frac{2}{8} = \frac{1}{4}. \qquad \square\]

This confirms that the ΔCov = 1/4 induction step observed computationally is
exact for odd-to-odd transitions. The constancy of the step size reflects the
fact that each increase of j by 2 adds exactly two new off-diagonal terms to
conv\_j, each contributing Cov = 1/8 by Lemma D.

\begin{center}\rule{0.5\linewidth}{0.5pt}\end{center}

\newpage

\subsection{5. Even-j Corrections}\label{even-j-corrections}

For even j, the convolution conv\_j contains a diagonal term g\_\{j/2\}
h\_\{j/2\} = g\_\{j/2\}², for which Lemma D does not apply. The decomposition
becomes:

\[\text{Cov}(\text{carry}_j, \text{conv}_j) = \underbrace{(j-2) \cdot \frac{1}{8}}_{\text{off-diagonal}} + \underbrace{\text{Cov}(\text{carry}_j,\; g_{j/2} h_{j/2})}_{\text{diagonal}}.\]

For the off-diagonal terms, Lemma D applies as before. The diagonal term
requires separate treatment. Conditioning on g\_\{j/2\} = 1 and h\_\{j/2\} = 1
fixes \emph{both} free bits at index j/2 simultaneously. At column j/2, neither
g\_\{j/2\} nor h\_\{j/2\} is available to flip the parity of total\_\{j/2\}, so
the Parity Lemma fails at that single column. The carry perturbation recursion
(Step 2 of the Lemma D proof) acquires a nonzero error at column j/2, which
propagates forward and decays geometrically but does not vanish by position j.
This introduces a correction:

\[\varepsilon_j = \text{Cov}(\text{carry}_j,\; g_{j/2} h_{j/2}) - \frac{1}{8}.\]

The correction ε\_j is strictly positive and decays exponentially with the gap
j/2 between the conditioning column and the target.

\textbf{Table of verified diagonal corrections.}

\begin{longtable}[]{@{}lllll@{}}
\toprule\noalign{}
j & ε\_j & Cov(carry\_j, conv\_j) & (j−1)/8 & Difference \\
\midrule\noalign{}
\endhead
\bottomrule\noalign{}
\endlastfoot
2 & 1/16 & 3/16 & 1/8 & 1/16 \\
4 & 1/32 & 13/32 & 3/8 & 1/32 \\
6 & 1/128 & 81/128 & 5/8 & 1/128 \\
8 & 1/256 & 225/256 & 7/8 & 1/256 \\
10 & 3/4096 & 4611/4096 & 9/8 & 3/4096 \\
12 & 11/32768 & 45067/32768 & 11/8 & 11/32768 \\
\end{longtable}

The corrections ε₂ = 1/16 and ε₄ = 1/32 have simple forms, but for j ≥ 10 the
numerators (3, 11, \ldots) do not follow a simple closed-form pattern. The
overall scaling ε\_j ≈ C · 2\^{}\{−j\} is consistent with the geometric
convergence rate of the perturbed carry chain after the Parity Lemma breaks at a
single column.

For all even j:

\[\text{Cov}(\text{carry}_j, \text{conv}_j) = \frac{j-1}{8} + \varepsilon_j,\]

so (j − 1)/8 serves as a lower bound with exponentially small correction. The
numerator pattern for j ≥ 10 is non-trivial (ε₁₀ = 3/4096, ε₁₂ = 11/32768),
reflecting the increasingly complex interaction between the diagonal
perturbation and the carry chain dynamics.

\begin{center}\rule{0.5\linewidth}{0.5pt}\end{center}

\newpage

\subsection{6. Computational Verification}\label{computational-verification}

All results have been verified by exact rational arithmetic, enumerating all
2\^{}\{2K\} bit configurations for K up to 12 (corresponding to 2\^{}\{24\} =
16,777,216 configurations).

\textbf{Enumeration method.} For a given K, we enumerate all 2\^{}\{2K\}
assignments of the K free bits in each of G and H (with g₀ = h₀ = 1 fixed),
compute the carry chain exactly using integer arithmetic, and accumulate the
required moments and cross-moments as exact rational numbers.

\textbf{Column and carry expectations.} For all j ≤ 12, E{[}conv\_j{]} = (j +
3)/4 and E{[}carry\_j{]} = (j − 1)/4 hold exactly.

\textbf{Off-diagonal universality.} At K = 8 (65,536 configurations per bit-pair
check), all off-diagonal covariances Cov(carry\_j, g\_i h\_\{j−i\}) for 1 ≤ i ≤
j − 1, i ≠ j − i, j ≤ 8, evaluate to exactly 1/8.

\textbf{Diagonal anomaly.} All diagonal covariances Cov(carry\_j, g\_\{j/2\}
h\_\{j/2\}) for even j ≤ 12 strictly exceed 1/8, with corrections matching the
table in §5.

\textbf{Odd-j formula.} For all odd j with 3 ≤ j ≤ 11, Cov(carry\_j, conv\_j) =
(j − 1)/8 exactly.

\textbf{ΔCov step.} The step Cov(carry\_\{j+2\}, conv\_\{j+2\}) − Cov(carry\_j,
conv\_j) = 1/4 holds exactly for all consecutive odd pairs (3, 5), (5, 7), (7,
9), (9, 11).

\textbf{Parity Lemma under conditioning.} The survival of the Parity Lemma under
conditioning on off-diagonal bit pairs (Step 1 of the Lemma D proof) has been
verified for all off-diagonal pairs with j ≤ 10: the conditional distribution of
total\_k mod 2 remains exactly Ber(1/2) at every column, confirming the
theoretical argument.

\textbf{Even-j diagonal corrections.} The diagonal covariance values
Cov(carry\_j, g\_\{j/2\} h\_\{j/2\}) for even j ≤ 12 all exceed 1/8, with the
excess matching the ε\_j values in §5 to full precision.

\begin{center}\rule{0.5\linewidth}{0.5pt}\end{center}

\newpage

\subsection{7. Discussion}\label{discussion}

\textbf{Independence from the trace anomaly conjecture.} The results in this
paper --- Lemmas A, B, D, Theorem 3, the Main Theorem, and the Corollary --- are
proved by elementary means (induction, conditioning, geometric sums). They do
not depend on the c₁ = π/18 conjecture, any spectral assumptions, or any
unproved distributional hypotheses about the carry chain. Every step relies only
on the independence structure of the bit model and the deterministic carry
recurrence.

\textbf{Connection to the non-Markovian correction.} The covariance
Cov(carry\_j, conv\_j) quantifies the non-Markovian correction δ\_NM in the
trace anomaly analysis. Specifically, the variance of carry\_j decomposes as:

\[\text{Var}(\text{carry}_{j+1}) = \frac{1}{4}\bigl(\text{Var}(\text{conv}_j) + \text{Var}(\text{carry}_j) + 2\,\text{Cov}(\text{carry}_j, \text{conv}_j) - \text{Var}(\text{total}_j \bmod 2)\bigr).\]

The exact evaluation of the covariance term removes one unknown from this
recursion, constraining the carry variance growth and thereby the trace anomaly.

\textbf{Odd vs.~even dichotomy.} The clean formula Cov = (j−1)/8 holds for all
odd j because the off-diagonal universality of Lemma D applies to every interior
term and no diagonal term exists. For even j, the single diagonal term
g\_\{j/2\} h\_\{j/2\} breaks the Parity Lemma under conditioning at column j/2,
introducing a correction that decays geometrically with the gap j/2 between the
conditioning point and the target column j. This odd/even dichotomy is a
structural feature of binary multiplication --- it arises from the parity of the
convolution length, not from any approximation in the proof method.

\textbf{Open question.} Whether the even-j corrections ε\_j admit a closed-form
expression remains open. The values through j = 12 are consistent with ε\_j =
f(j) · 2\^{}\{−j\} for some slowly-varying function f, but the numerator
sequence (1, 1, 1, 1, 3, 11, \ldots) does not match any standard combinatorial
sequence.

\textbf{The perturbation method.} The proof of Lemma D introduces a general
technique: condition on a bit pair, track the perturbation to column
expectations via an exact linear recursion (enabled by the Parity Lemma), and
evaluate a weighted geometric sum. This method extends naturally to higher-order
covariances and to other bases, suggesting that analogues of the universal
off-diagonal covariance may hold in broader settings.

\textbf{Companion papers.} The spectral theory of the carry operator is
developed in {[}A{]} (\emph{Spectral Theory of Carries in Positional
Multiplication}). The trace anomaly and its connection to the c₁ conjecture are
analyzed in {[}E{]} (\emph{The Trace Anomaly of Binary Multiplication}). The
present paper provides the exact covariance inputs required by both analyses. In
particular, the Main Theorem supplies the Cov(carry\_j, conv\_j) term that
governs the non-Markovian correction in the variance recursion of {[}E{]}, and
the Parity Lemma provides the key parity input for the spectral analysis of
{[}A{]}.

\begin{center}\rule{0.5\linewidth}{0.5pt}\end{center}

\newpage

\subsection{8. References}\label{references}

\begin{enumerate}
\def\labelenumi{\arabic{enumi}.}
\item
  P. Diaconis and J. Fulman, ``Carries, shuffling, and symmetric functions,''
  \emph{Advances in Applied Mathematics}, vol.~43, no. 2, pp.~176--196, 2009.
\item
  J. M. Holte, ``Carries, combinatorics, and an amazing matrix,'' \emph{The
  American Mathematical Monthly}, vol.~104, no. 2, pp.~138--149, 1997.
\item
  D. E. Knuth, \emph{The Art of Computer Programming, Volume 2: Seminumerical
  Algorithms}, 3rd ed.~Addison-Wesley, 1997.
\item
  {[}A{]} Companion paper: \emph{Spectral Theory of Carries in Positional
  Multiplication}, this series. doi:10.5281/zenodo.18895593
\item
  {[}E{]} Companion paper: \emph{The Trace Anomaly of Binary Multiplication},
  this series. doi:10.5281/zenodo.18895604
\item
  {[}P2{]} Companion paper: \emph{The Sector Ratio in Binary Multiplication:
  From Markov Failure to Transcendence}, this series.
  doi:10.5281/zenodo.18895615
\end{enumerate}

\end{document}
